\section{Pairing down to a sparse boundary}\label{sec:early-phase}

We now construct a partial pairing with $l_1$ unresolved stubs. Throughout
this phase, both signed precisions remain positive and their normalized
inverses remain close, entry by entry, to the truncated path matrix $P$.
The structural estimates of the preceding section then leave fewer than
$b$ bad stubs for the separation step to remove.

All conditional expectations in this section are over a uniform ordered
pair of distinct unresolved stubs whose vertices have distance greater than
\(g_0\). A conditional expectation for a surviving output additionally
excludes its one or two stub labels from the selected pair. Global vertex
outputs do not impose this exclusion. The notation \(E_{ij}\) always uses
the vectors attached to the current labels; thus a stub output is centered
using the current unresolved set.

\begin{proposition}[Early pairing]\label{prop:early-phase}
For the fixed parameters of the construction and all sufficiently large
\(n\), the simultaneous discounted-account rule can be run from the empty
pairing down to \(l_1\) unresolved stubs. It retains the structural estimates
\eqref{eq:original-3}--\eqref{eq:original-7}, positivity of every tracked
precision, and, at every tracked point,
\[
 |L_1|\le n^{.075a}/\alpha_*,\qquad
 |L_0|\le n^{.105a}/\alpha_*^2,\qquad
 \max_{i,j}|E_{ij}|\le n^{.33a}/\sqrt{l_0}.
\]
The low-square estimates \eqref{eq:original-23} hold with any prescribed
sufficiently small fixed loss exponent, for any fixed near radius needed
in the proof. At every primary point,
\begin{equation}
 \mathcal T \ \ge\ n^{-\theta/2}(n^{-1}+\alpha/l)^{-1}
\label{eq:original-13}
\end{equation}
after a fixed constant is allowed in the right-hand side.
\end{proposition}

At every exposure, the structural accounts, the trace accounts, the
squared entry sums, and the high entry moments enter the same discounted
total. The proofs below establish their one-step bounds under the common
current-state caps. The squared sums make mixed terms in the high-moment
expansion summably small. A separate weighting of pair and diagonal
moments controls the remaining terms with a drift coefficient independent
of the moment order. This independence permits extraction of a
simultaneous entry bound after a polynomial loss in the discount.

At the empty state, \(W=P=I|_{{\bf1}^\perp}\),
\(L_1=1-d\), \(L_0=0\), and \(\mathcal T=d(n-1)\).
Write
\[
 B_e=n^{.31a}/\sqrt{l_0},\qquad d_m=n^{.33a}/\sqrt{l_0}.
\]
We maintain \(\max|E_{ij}|\le d_m\) together with the squared-sum bounds
stated below; in particular,
\(T_S:=\sum_{x\in S} E_{xx}^2\le n^\eta\), with centered stub indices.
The deterministic implications \eqref{eq:original-8}--\eqref{eq:original-12}
are used only at the current state. There, every valid pair has small
\(D_p\), so its next precisions are positive and the update
\eqref{eq:original-1} is defined. The weighted choice and extraction then
give the displayed caps at the selected next state. This closes the
induction without assuming any candidate's next-state estimates.

\subsection{Trace estimates}

We first control the two trace errors $L_0,L_1$. We also bound
$\mathcal T$ away from zero; this lower bound will provide a strict
margin in the boundary-completion argument. The signed cancellations
needed for these three estimates are collected in the next lemma.

Put \(D=\Gamma-\Pi_l\), a centered full stub matrix with pair block
\(D_p\), and put \(J_{\rm all}=\Gamma Q_{\rm all}\).
In this subsection, \(K,Q,J\) denote the selected pair blocks
of \(K_{\rm all},Q_{\rm all},J_{\rm all}\), respectively.
At the current state,
\begin{equation}
 K_{xx}\le C,\quad Q_{xx}\le C/\alpha,\quad
 \sum_x D_{xx}^2\le n^\eta,\quad \sum_y D_{xy}^2\le C,
\label{eq:original-14}
\end{equation}
by \eqref{eq:original-4}, the squared-sum bound, and
\eqref{eq:original-10}. On valid pairs, \(D_p\) has power-small entries;
on invalid pairs its entries are still bounded. Fixed constants are
included in the expectation-error bounds below. We retain the exact
ordered-distinct denominators in the identities that supply the drift.

\begin{lemma}[Uniform pair contractions]\label{lem:early-pair-contractions}
Let \(A,C\) be centered \(l\)-by-\(l\) matrices, with \(A\) symmetric,
and let a subscript \(p\) denote their \(\{u,v\}\) blocks. For a full
uniform ordered distinct pair,
\[
\begin{split}
 \mathbb E\operatorname{Tr}(H_0A_p)&=0,\\
 \mathbb E\operatorname{Tr}(BA_p)&=\frac{2\operatorname{Tr}A}{l-1},\\
 \mathbb E\operatorname{Tr}(FA_pFC_p)
 &=\frac{2}{l(l-1)}\left(
       \operatorname{Tr}A\operatorname{Tr}C+
       \operatorname{Tr}(AC)-2\sum_x A_{xx}C_{xx}\right).
\end{split}
\]
The last identity remains valid when \(C\) is not symmetric. For a
nonnegative summand, restriction to valid pairs or to survival costs a
factor \(1+O((D_g+1)/l)\). For a signed summand one must also subtract
the omitted pairs before applying this factor.
\end{lemma}
\begin{proof}
Centering gives \(\sum_{u\ne v}A_{uv}=-\operatorname{Tr}A\).
Insert \(B=I+J_2/(l-2)\) and
\(H_0=-\sigma h(F+J_2/(l-2))\) to obtain the first two
identities. Expansion of the last trace gives
\(A_{vv}C_{uu}+A_{uu}C_{vv}+A_{uv}(C_{uv}+C_{vu})\).
Summing its three terms proves the formula. Each stub has at most
\(D_g\) forbidden partners, and survival excludes at most two stubs;
the assertion about restriction follows by counting the available
ordered pairs.
\end{proof}

\begin{lemma}[Trace transition bounds]\label{lem:early-trace-transitions}
Under the current bootstrap bounds, the conditional means and second
moments of the trace errors satisfy
\eqref{eq:original-16}--\eqref{eq:original-18} at all tracked points.
At primary points, \(\mathcal T\) has the lower relative drift
\eqref{eq:original-20}, the second-moment bound
\eqref{eq:early-T-variance}, and relative step size
\(\lvert\Delta\mathcal T\rvert/\mathcal T=o(1)\).
All error constants are uniform in the current state.
\end{lemma}
\begin{proof}

We first compute the two centered trace increments. We then obtain the
three estimates on \(\mathcal T\) needed to bound its reciprocal:
the lower drift, second moment, and maximum relative increment. The terms
containing the possibly large matrix \(J\) require separate estimates.

\emph{The centered trace errors.}

From \(\mathcal D'=\mathcal D-U B U^T\),
\[
 \Delta F_1=-\operatorname{Tr}B G-\operatorname{Tr}H K+\operatorname{Tr}B G H G .
\]
Expand \(H=H_0-H_0 D_p H_0+\cdots\). In the increment of \(L_1\),
the constant \(-2\) cancels the change in its baseline \(l-1\). This gives
\begin{equation}
 \Delta L_1=
 2\operatorname{Tr}H_0 D_p-(1+h^2)\operatorname{Tr}B D_p-\operatorname{Tr}H_0 K
 +\operatorname{Tr}H_0 D_p H_0 K+O(\|D_p\|^2).
\label{eq:original-15}
\end{equation}
Here \(H_0 G_0 H_0=h^2 B\). Under full uniform distinct sampling, the
linear \(H_0\) contractions vanish and
\(\mathbb E\operatorname{Tr}B D_p=2L_1/(l-1)\).
The swap-swap term has diagonal contribution
\(2 h^2 L_1\operatorname{Tr}K_{\rm all}/[l(l-1)]\), with same-index
corrections \(O(1/l)\). Its remaining cross contribution is bounded by
\(C(\operatorname{Tr}\Gamma^3+l)/l^2\); the \(O(1/l)\) changes in swap
entries also cost \(O(1/l)\).

To restrict these polynomial terms to valid pairs, subtract an absolute
error \(O(D_g/l)\) and divide by \(1-O(D_g/l)\). The remaining inverse
series has expectation \(O(n^\eta/l)\) by \eqref{eq:original-14}.
Consequently,
\begin{equation}
 \mathbb E\Delta L_1=a_l L_1+
 O\!\left((n^\eta+D_g)/l+n^c/(l\alpha)+n^{3v_*}/(l^2\alpha^3)\right),
 \quad
 -C/l\le a_l\le n^{-c_3}/l
\label{eq:original-16}
\end{equation}
for some \(c_3>0\). Before conditioning, the coefficient is
\(\frac{2}{l-1}[h^2(\operatorname{Tr}K_{\rm all}/l-1)-1]\).
Its stated bounds follow from the averaged upper bound in
\eqref{eq:original-10} and \(\|\xi_x\|^2=1+O(1/l)\).
The positive normalization for valid pairs preserves these upper and
lower bounds. Also,
\begin{equation}
 \mathbb E(\Delta L_1)^2\le C\big(n^{2c}/(l\alpha^2)+n^{4v_*}/(l^2\alpha^4)+n^\eta/l\big).
\label{eq:original-17}
\end{equation}
Indeed, the ideal \(K_{uv}\) term uses \eqref{eq:original-12}, while
diagonal same-end ideal terms have coefficient \(O(1/l)\).
All other terms use \eqref{eq:original-14}.

For the global trace, \(\Delta L_0=-\operatorname{Tr}H Q\).
The ideal centered term has zero full mean. The term with a single
\(D_p\) between pure swaps has diagonal main contribution at most
\(C|L_1|/(l\alpha)\). Its cross contribution is bounded using
\(\operatorname{Tr}(\Gamma+\Pi_l)Q_{\rm all}\), together with the
same-index corrections. The remainder is \(O(n^\eta/(l\alpha))\).
Thus, on valid choices,
\begin{equation}
\begin{split}
 |\mathbb E\Delta L_0|&\le C\big((|L_1|+n^\eta+D_g)/(l\alpha)
       +n^c/(l\alpha^2)+n^{3v_*}/(l^2\alpha^4)\big),\\
 \mathbb E(\Delta L_0)^2&\le C\big(n^{2c}/(l\alpha^4)+n^{4v_*}/(l^2\alpha^6)
                               +n^\eta/(l\alpha^2)\big).
\end{split}
\label{eq:original-18}
\end{equation}
The contractions used here can be read directly from
Lemma~\ref{lem:early-pair-contractions}: for \(K_{\rm all}=\Gamma^2\),
\(\operatorname{Tr}(DK_{\rm all})=
\operatorname{Tr}(\Gamma^3)-\operatorname{Tr}(\Gamma^2)\), and for
\(Q_{\rm all}\succeq0\),
\(\lvert\operatorname{Tr}(DQ_{\rm all})\rvert
\le\operatorname{Tr}(\Gamma Q_{\rm all})+\mathcal T\).
The diagonal exclusions are bounded using
\(K_{xx}\le C\), \(Q_{xx}\le C/\alpha\), and
\(\sum_x|D_{xx}|\le n^{\eta/2}\sqrt l\).
Terms with two or more \(D_p\) factors have absolute value at most
\(C\|D_p\|^2\), respectively \(C\|D_p\|^2/\alpha\), because
the remaining inverse series is geometric. Their expectations are
bounded by \eqref{eq:original-14}. The ideal variances are
\(C\operatorname{Tr}(K_{\rm all}^2)/l^2\) and
\(C\operatorname{Tr}(Q_{\rm all}^2)/l^2\), respectively; applying
\eqref{eq:original-12} gives exactly the first two summands of
\eqref{eq:original-17} and \eqref{eq:original-18}.

\emph{Lower drift of \(\mathcal T\).}
To prove \eqref{eq:original-13}, we will control the reciprocal statistic.
We first bound \(\mathbb E\Delta\mathcal T\) from below.
The exact update of $\mathcal T$ is
\begin{equation}
 \Delta\mathcal T=
 -\operatorname{Tr}\big[B(I-GH)Q(I-HG)\big]-2\operatorname{Tr}H J
       +\operatorname{Tr}H Q H K .
\label{eq:original-19}
\end{equation}
Using \(G_0,H_0\) in the outer terms, their full-pair mean divided by \(\mathcal T\) is
\[
 \frac{2}{l-1}\,[h^2(\operatorname{Tr}K_{\rm all}/l-1)-1]
       + O((1+\|\Gamma\|^2)/l^2).
\]
For the first outer term, expand \(I-G_0H_0=I+\sigma h F\).
The centered \(F B\) term averages to zero, and the other contribution is
\(-2(1+h^2)/(l-1)\). For the last term, use the same opposite-diagonal
swap average. Its cross trace satisfies
\(\operatorname{Tr}K_{\rm all} Q_{\rm all}\le\|\Gamma\|^2\mathcal T\);
the diagonal exclusions and \(O(1/l)\) swap errors use \(K_{xx}\le C\).

At primaries, \eqref{eq:original-11} therefore bounds this relative drift
below by \(-C\alpha\mathcal T/l^2 -n^{-c_4}/l\). Pair smallness and
\eqref{eq:original-14} bound the change from the ideal outer terms, on
valid pairs, by \(n^{-c_4}\operatorname{Tr}Q\), decreasing
\(c_4>0\) if necessary. The invalid ideal contributions have average
\(O(D_g\mathcal T/l^2)\).

For the \(J\) term, the ideal contraction \(H_0J\) has zero full
mean. In the term with one deviation, \(H_0D_pH_0J\), diagonal
factorization uses
\(\operatorname{Tr}J_{\rm all}=\operatorname{Tr}\Gamma Q_{\rm all}\),
and the cross trace uses \(D=\Gamma-\Pi_l\). The swap average is
therefore bounded by
\[
 C\bigl(|L_1|\|\Gamma\|+\|\Gamma\|^2+\|\Gamma\|\bigr)
       \mathcal T/l^2,
\]
apart from same-index exclusions, invalid pairs, and the
\(O(1/l)\) changes in the swap coefficients. This main bound is
power-small relative to \(\mathcal T/l\), since
\(\lvert L_1\rvert\le n^{.075a}/\alpha_*\) and
\(n^{2v_*}/(l\alpha_*^2)\ll1\).

To handle the exclusions and coefficient changes, including the
invalid-pair part of the ideal \(J\) contraction, use
\[
 |J_{xy}|\le C\sqrt{(Q_{\rm all}^2)_{yy}},\quad
 \sum_y\sqrt{(Q_{\rm all}^2)_{yy}}\le C n^{v_*}\alpha^{-3/2}\sqrt{l\mathcal T}.
\]
Their mean, divided by \(\mathcal T/l\), is at most
\[
 \frac{C(D_g+1)n^{v_*}}{\alpha^{3/2}\sqrt{l\mathcal T}},
\]
which is power-small under the current lower bound
\eqref{eq:original-13}. Finally, terms with at least two deviations
use \(\mathbb E_{\rm valid}\|D_p\|^2\le Cn^\eta/l\) and
\(\max|J_{xy}|\le Cn^{v_*}/\alpha^2\); their contribution is
also power-small relative to \(\mathcal T/l\). Conditioning on
valid pairs preserves the resulting lower drift bound. We have proved
\begin{equation}
 \mathbb E\Delta\mathcal T/\mathcal T \ \ge\ -C\alpha\mathcal T/l^2-n^{-c_4}/l
\label{eq:original-20}
\end{equation}
with some \(c_4>0\), for sufficiently small \(\theta,\eta\).

\emph{Second moment and maximum increment of \(\mathcal T\).}
In \eqref{eq:original-19}, the outer terms are bounded by
\(C\operatorname{Tr}Q\), giving second moment
\(C\mathcal T/(\alpha l)\). The off-diagonal \(J\) terms use
\(\|\Gamma Q_{\rm all}\|_F^2\le
Cn^{4v_*}\mathcal T/\alpha^5\).
For a diagonal \(J_{uu}\), the coefficient \(H_{uu}\) has
conditional squared average at most \(Cn^\eta/l\) over a uniform
valid partner. Indeed, expansion around \(H_0\) gives
\begin{equation}
 |H_{uu}|\le C(1/l+|D_{vv}|+|D_{uv}|)
\label{eq:original-21}
\end{equation}
on valid pairs: the first pure swap from \(u\) touches \(v\),
all other factors are uniformly small, and a non-pure initial swap
costs \(O(1/l)\). The conditional claim follows from
\eqref{eq:original-14} and \(D_g\ll l\). Also, the norm bound
times the trace gives
\[
 \sum_u J_{uu}^2\le C\operatorname{Tr}Q_{\rm all}^2
       \le Cn^{2v_*}\mathcal T/\alpha^3.
\]
Consequently,
\begin{equation}
 \mathbb E(\Delta\mathcal T)^2\le C\left(\frac{\mathcal T}{\alpha l}+
       \frac{n^{4v_*}\mathcal T}{\alpha^5 l^2}
       +\frac{n^{2v_*+\eta}\mathcal T}{\alpha^3 l^2}\right).
\label{eq:early-T-variance}
\end{equation}
Individually, \(\lvert\Delta\mathcal T\rvert\le
Cn^{v_*}/\alpha^2\ll\mathcal T\) under
\eqref{eq:original-13}.
\end{proof}

\begin{lemma}[Trace control by weighted choice]\label{lem:early-trace-accounts}
The trace-square accounts and the inverse-\(\mathcal T\) account have
bounded discount loss under the bootstrap bounds and imply the trace
conclusions of Proposition~\ref{prop:early-phase} with a positive power
of slack.
\end{lemma}
\begin{proof}
\emph{Trace-square accounts.}
Track \(1+L_i^2/S_i^2\), where
\(S_1=n^{.065a}/\alpha_*\) and
\(S_0=n^{.095a}/\alpha_*^2\).
In \eqref{eq:original-16}, the coefficient contribution to square drift
is at most \(2n^{-c_3}L_1^2/l\). Use
\(2|L|/S\le1+L^2/S^2\) for each additive mean error.
Since \(\sum l^{-1}=O(\log n)\),
\(\sum l^{-2}=O(l_0^{-1})\), and
\(\alpha_*^2l_0\asymp n^{1.15a}\), the normalized mean and variance
budgets for \(L_1\) are bounded by fixed multiples of
\[
 (n^{\eta-.065a}\alpha_*+n^{-.063a}\alpha_*+
   n^{-.035a})\log n+n^{-.840a}
 \quad\hbox{and}\quad
 n^{-.070a}\log n+n^{-.780a}+n^{\eta-.130a}\alpha_*^2\log n.
\]
Both tend to zero. The \(L_0\) bounds follow from
\eqref{eq:original-18}: the term involving \(L_1\) loses at most
\(n^{-.020a}\log n\), and every other normalized term has a larger
strict margin. This proves the two trace-square assertions.

\emph{The reciprocal account.}
Divide \eqref{eq:early-T-variance} by \(\mathcal T^2\) and sum
using the current lower bound \eqref{eq:original-13}. The result is
power-small; the particularly relevant bounds are
\(n^{4v_*+\theta}(1/(n\alpha^5 l_0)+1/(\alpha^4 l_0^2))\ll1\).
The maximum relative increment from
Lemma~\ref{lem:early-trace-transitions} permits Taylor expansion of
the reciprocal. Together with \eqref{eq:original-20}, this yields
\[
 \mathbb E(1/\mathcal T')\le (1+u_l)/\mathcal T + C\alpha/l^2,
 \qquad \sum u_l=o(1)
\]
with deterministic nonnegative upper budgets under the caps.
For \(f_l=1/n+\alpha/l\), use the account
\(1+n^{-\theta/4}/(\mathcal T f_l)\). The scaled additive errors
\(Cn^{-\theta/4}\alpha/(l^2f_{l-2})\) sum to \(o(1)\).
This proves \eqref{eq:original-13} with a positive power of slack.

\end{proof}

\subsection{Entry increments and their conditional means}

The trace bounds do not alone control each matrix entry. We therefore
compare the exact inverse update with the truncated path update. The
resulting expansion has two kinds of terms whose signs must be retained;
all other terms can be estimated by absolute path bounds. The following
classification records every possible position of an error factor.

In the update expansion below, entries are first evaluated on the old
output vectors. Recentring the surviving stub vectors is dealt with explicitly
at the end of Lemma~\ref{lem:early-entry-mean}.
For a matrix \(X\) on vertex space put
\[
 \mathcal Q(X)=XU H_X U^T X,\qquad
 H_X=H_0\sum_{m\ge0}[-(U^TXU-G_0)H_0]^m.
\]
For \(X=W\) or \(P\) these series converge numerically on every valid
pair. The formula for \(\mathcal Q(P)\) also has its coefficientwise
meaning from \eqref{eq:original-1}. Set
\[
 T=U^TPU-G_0,\qquad D_E=U^TEU,
 \qquad \mathcal R_P=[\mathcal Q(P)]_{\ge L}.
\]
We call an entry of \(T\) a topology factor: it measures the path
reference's deviation from the pair comparison block \(G_0\).
A diagonal topology factor at \(x\) is bounded by \(j_x\), and
an off-diagonal one by \(b_{uv}\). We also call diagonal topology
factors local factors. Entries of \(D_E\) are intermediate error
factors; errors in the two outer entries will be kept separate.
The exact increment on the old output vectors is
\[
 \Delta E=-\{\mathcal Q(W)-\mathcal Q(P)\}-\mathcal R_P.
\]
This identity fixes the sign of the formal tail. We retain signs for the
terms with no intermediate factor and for the terms with one
intermediate error and no topology factor. The proof below identifies
their cancellations explicitly.

For an output label \(i\), let
\(r_i(x)=A_i(x)+|e_i(x)|\). The hypotheses give
\[
 \sum_{x\in S}r_i(x)^2\le C,\quad
 \max_x r_i(x)\le C,\quad
 \sum_x j_x^2\le n^\eta,\quad
 \max_u\sum_v b_{uv}^2\le C,
 \quad Q_a:=n^\eta q^{L/2}.
\]
Increasing \(n^\eta\) by an arbitrarily small exponent if necessary,
\(\sum_xj_x\le Q_a\) and
\(\max_u\sum_v b_{uv}\le Q_a\), by
\eqref{eq:original-4} and \eqref{eq:path-row-sum}, respectively.
All matrices of majorants on stub labels can be taken symmetric.

\begin{lemma}[Index-word classification]\label{lem:early-index-words}
Expand every occurrence of \(H_0\) as
\(-\sigma h F-\sigma hJ_2/(l-2)\), every intermediate factor as
\(T+D_E\), and each outer factor as \(P+E\). A resulting scalar word
has two outer endpoint indices \(s,t\in\{u,v\}\), alternating
swap and intermediate factors, and one of the following descriptions.
\begin{enumerate}
\item It contains a non-pure swap \(-\sigma hJ_2/(l-2)\).
\item All swaps are pure and it contains an off-diagonal \(T\) factor,
      bounded by \(b_{uv}\).
\item All swaps are pure, all \(T\) factors are diagonal, and at least
      two intermediate factors are \(D_E\).
\item All swaps are pure, all \(T\) factors are diagonal, there is
      at least one intermediate factor, and there is at most one
      \(D_E\) factor. If \(s=t\), either there is a
      diagonal \(T\) factor at \(\bar s\), or the word has exactly
      one intermediate factor, namely \((D_E)_{\bar s\bar s}\).
      If \(s\ne t\) and there is a \(T\) factor, there is at least
      one local factor \(j_u\) or \(j_v\). If there is no \(D_E\)
      factor, both local factors occur.
\item There are no intermediate factors. Its pure-swap portion has
      \(s\ne t\).
\end{enumerate}
Here \(\bar u=v\) and \(\bar v=u\). The list is exhaustive; the cases
are made disjoint by testing them in the displayed order. With only
one intermediate \(D_E\) and no \(T\), the exceptional same-end
word in item 4 is precisely
\(h^2 x_i(s)x_j(s)e_{\bar s}(\bar s)\), where \(x_i,x_j\)
are the chosen outer arrays, each a \(p\) or \(e\) row.
All estimates obtained by retaining at most a fixed number of marked
factors in these cases can be summed over the remaining words with a
constant independent of \(n\).
\end{lemma}
\begin{proof}
After item 1 is excluded, a swap changes an endpoint to its opposite.
An intermediate diagonal \(T\) leaves that endpoint unchanged. An
intermediate off-diagonal \(T\) puts the word in item 2. Exclude
items 2 and 3. In a same-end word, the index immediately after the
first swap and immediately before the last swap is \(\bar s\).
If a local factor at \(\bar s\) never occurs, both these visits must
be incident to the sole \(D_E\) factor. A swap between distinct
intermediate positions changes the index, so a second intermediate
position would force either another \(D_E\) or a local factor at
\(\bar s\). Hence the sole intermediate is
\((D_E)_{\bar s\bar s}\). In a distinct-end word consisting only of
local factors, the number of intermediates is positive and even;
the indices alternate, so both endpoints occur. These observations
prove items 4 and 5 and their exhaustiveness.

For convergence, let \(\varepsilon_n\) bound the absolute entry sum
of an intermediate \(T+D_E\) block on a valid pair.
Equations \eqref{eq:original-5} and the entry bootstrap imply
\(\varepsilon_n\to0\). There are at most \(4^{m+1}\) endpoint
index strings with \(m\) intermediates. If at most \(b\) positions
are marked, summation of the unmarked positions is dominated by
\[
 C_b\sum_{m\ge b}(m+1)^b(C\varepsilon_n)^{m-b},
\]
which is bounded for large \(n\). A non-pure swap at a marked
position contributes \(C/l\). The coefficientwise absolute
majorants of the tail have the same convergence after its degree
restriction is discarded. This proves the final assertion and also
justifies termwise expectation and the use of Minkowski's inequality
for second moments.
\end{proof}

The elementary averaging bounds behind the classification will be used
repeatedly. For nonnegative arrays \(a,c\) with bounded squared sums,
full independent uniform endpoints satisfy
\[
\begin{array}{c|c|c}
\text{retained factors}&\text{absolute first moment}&
                       \text{second moment}\\ \hline
 a(u)c(v)&C/l&C/l^2\\
 a(u)c(u)/l&C/l^2&C/l^3\\
 a(u)c(u)j_v& C Q_a/l^2&C n^\eta/l^2\\
 a(u)c(u)b_{uv}& C Q_a/l^2&C/l^2\\
 a(u)c(v)j_u& C n^{\eta/2}/l^{3/2}&C n^\eta/l^2\\
 a(u)c(v)b_{uv}&C Q_a/l^2&C/l^2\\
 d_m^2a(u)c(u)&C d_m^2/l&C d_m^4/l.
\end{array}
\]
For example, the first-moment cross-topology bound follows from
\(\|b\|\le\max_u\sum_vb_{uv}\le Q_a\); the same-end version
uses its row-sum bound before summing \(a(u)c(u)\). For second
moments of a cross-topology term, sum \(b_{uv}^2\) over its free
endpoint and use the bounded suprema of the outer arrays. For a local
factor, use \(\sum_vj_v^2\le n^\eta\), or Cauchy--Schwarz for its
first moment. These arguments prove every table entry. Restricting
to distinct, valid, surviving endpoints costs a bounded density
factor. Consequently the table also applies under the conditional
laws of this section.

\begin{lemma}[Signed entry mean]\label{lem:early-entry-mean}
For every existing output pair, conditionally on its survival,
\begin{equation}
 \big|\mathbb E\Delta E_{ij}\big|\le \mu_l,
 \qquad
 \mu_l=C\left[
 \frac{n^{v_*}/\alpha_*+n^\eta q^{L/2}+D_g}{l^2}
 +\frac{n^\eta}{l^{3/2}}+\frac{d_m^2}{l}\right].
\label{eq:original-22}
\end{equation}
This includes recentered stub outputs.
\end{lemma}
\begin{proof}
First consider words without intermediate factors. Their exact
\(H_0\) contraction is zero for a full ordered distinct pair by
\eqref{eq:original-2}; the same statement holds coefficientwise
for the pure-\(P\) tail. Omitting at most two labels costs
\(C/l^2\), using their bounded entries and the squared row bounds
in \eqref{eq:original-2}. If \(\mathcal B\) is the symmetric
forbidden-pair indicator, then \(\|\mathcal B\|\le D_g\).
Thus an opposite-outer product on forbidden pairs has sum at most
\(D_g\|r_i\|_2\|r_j\|_2\le CD_g\). A same-end product has
at most \(D_g\) partners and an extra \(1/l\) from its
non-pure swap. After division by \(l(l-1)\), the forbidden-pair
cost is \(CD_g/l^2\).

Next retain exactly one intermediate \(D_E\), no topology factor,
pure swaps, and arbitrary \(W\) outers. This groups all outer
choices without losing their signs. Write
\[
 Y_E=\operatorname{Tr}(\Xi^TE\Xi),\qquad
 |Y_E|\le |L_1|+Q_a,\qquad
 \|\Xi^TE\Xi\|\le Cn^{v_*}/\alpha+C(1+L)^3.
\]
For same outer endpoint \(u\), the numerator sum is
\[
 \sum_{u\ne v}w_i(u)w_j(u)e_v(v)
 =Y_E\sum_u w_i(u)w_j(u)
   -\sum_u w_i(u)w_j(u)e_u(u),
\]
whose magnitude is \(C(1+|Y_E|)\). For distinct outer endpoints
it is a bilinear contraction of \(\Xi^TE\Xi\) with two bounded
\(\ell^2\) rows, after a diagonal exclusion, and is bounded by
\(C(1+\|\Xi^TE\Xi\|)\). On a surviving set \(I\) obtained
by deleting at most two output labels, the same-end numerator is
\[
 \left(\sum_{v\in I}e_v(v)\right)
 \left(\sum_{u\in I}w_i(u)w_j(u)\right)
 -\sum_{u\in I}w_i(u)w_j(u)e_u(u),
\]
and is still bounded by \(C(1+|Y_E|)\): the first diagonal sum
changes by \(O(d_m)\), and the second factor is bounded by the
two row norms. For distinct outer endpoints, deleting a bounded
number of rows or columns changes the bilinear numerator by
\(C\). The required bounded \(\ell^2\) rows of \(\Xi^TE\Xi\)
follow from \(E=W-P\) and the stub-row square bounds for \(W\)
and \(P\) in \eqref{eq:original-10} and \eqref{eq:original-4}.
Forbidden-pair subtraction costs \(CD_g\) by the preceding
operator-norm argument; the extra intermediate entries are bounded.
These estimates give the \(n^{v_*}/\alpha_*\) contribution in
\eqref{eq:original-22}.

Every remaining word is covered by Lemma~\ref{lem:early-index-words}.
A non-pure swap gives \(C/l^2\) in first moment. A cross-topology
word gives \(CQ_a/l^2\). If there is no cross topology and at most
one intermediate \(D_E\) factor, a same-end topology word contains
\(j_{\bar s}\), giving \(CQ_a/l^2\), and a distinct-end topology
word contains a local factor, giving \(Cn^\eta/l^{3/2}\).
A word with at least two intermediate \(D_E\) factors gives
\(Cd_m^2/l\). These statements apply also when an outer is \(E\),
since it is bounded by \(r_i\). They also apply to the coefficient
tail by absolute majorization. The geometric summation in the
classification lemma gives the same bounds for the entire series.

Finally, for a surviving stub,
\(\xi'_x=\xi_x+(\xi_u+\xi_v)/(l-2)\).
The old-axis entries of the new \(E\) have sup \(Cd_m\): this
follows from the difference series, in which every numerical word
contains an \(E\), and from \eqref{eq:original-6}. For a fixed
other output \(i\),
\(\mathbb E|E'_{ui}|+\mathbb E|E'_{vi}|\le C/\sqrt l\).
Indeed the old entry has this bound by its squared row sum. In every
update word the selected-axis outer is bounded and the other outer
is at \(u\) or \(v\); averaging that outer uses
\(\sum_xr_i(x)\le\sqrt l\|r_i\|_2\).
The formal tail is bounded coefficientwise by the same argument.
Thus the two first-order recentering terms have mean
\(Cl^{-3/2}\), and the term changing both axes has mean
\(Cd_m/l^2\). Their pointwise total is \(Cd_m/l\).
This proves the stated bound.
\end{proof}

\subsection{Low square sums}

The row-square estimates for \(W\) and \(P\) already control averages
of the squared entries of \(E\) in a row. We now bound the same errors
on diagonals and among pairs at a fixed distance. The exceptional same-end
word in Lemma~\ref{lem:early-index-words} copies a selected endpoint's
diagonal error to surviving outputs. For stub diagonals, deletion of the
two selected labels will outweigh this contribution. For nearby pairs,
we must also count those brought close by the new edge. We first bound
the reference entries that arise in this second calculation.

\begin{lemma}[Aligned short-path residual]\label{lem:early-near-reference}
Fix an integer \(k\). For \(i\in X\), \(u\in S\) at actual
vertex distance at most \(k\), remove from the actual part of
\(p_i(u)\) the monomial of the unique short path. The squared sum
of the remaining coefficientwise absolute majorants over all these
pairs is at most
\[
 n^\eta+C(1+L)^4q^L/l.
\]
Here \(X=V\) or \(S\), and \(\eta\) can be any prescribed positive
constant after the structural losses are reduced.
\end{lemma}
\begin{proof}
For large \(n\), \(g_0>2k\), so the short path is unique. Expand
the square of the remaining actual paths. A summand consists of that
short path and two nonbacktracking alternatives of lengths
\(a_1,a_2<L\). Its union has a cycle; the alternative equal to the
short path was removed. If the excess is at least two, the
structural diagram estimate gives, after multiplication by
\(h^{a_1+a_2}\), at most
\(n^{-1+o(1)}q^{(a_1+a_2)/2+k}\). Summing the logarithmically many
lengths is still at most \(n^{-1+o(1)}q^{L+k}=o(1)\).

For excess one the union is unicyclic. Suppress its forced tree tails
at the two endpoints. A nonbacktracking walk can only follow one of
the two cycle arcs, possibly with complete laps. If the short path
uses one arc, every alternative distinct from it covers the other
arc or completes a lap. If the short path stays in an attached tree,
every alternative reaches the cycle and completes a lap. Consequently
all edges outside the short path belong to each alternative, and
\(e_H\le\min(a_1,a_2)+k\). The weighted embedding bound is therefore
\[
 n^{o(1)}q^{e_H-(a_1+a_2)/2}
 \le n^{o(1)}q^k.
\]
The number of shapes, length choices, and walk choices is absorbed
by an arbitrarily small fixed power. Bounded stub multiplicities
change only constants. Finally there are \(O_k(l)\) near pairs,
and the per-entry uniform and centering error is
\(C(1+L)^2q^{L/2}/l\). Its square sum has the displayed size.
\end{proof}

\begin{lemma}[Control of squared entry sums]\label{lem:early-low-squares}
For any fixed radius \(R_c\), the simultaneous accounts retain
\begin{equation}
\begin{split}
 &T_S=\sum_{x\in S}e_x(x)^2\le n^\eta,
 \qquad T_V=\sum_{i\in V}E_{ii}^2\le n^\eta n/l_0,\\
 &T_{XS,k}=\sum_{\substack{i\in X,\ x\in S\\
    \operatorname{dist}(\operatorname{vtx}i,\operatorname{vtx}x)\le k}}
           e_i(x)^2\le n^\eta,
 \qquad X\in\{V,S\},\quad 0\le k\le R_c .
\end{split}
\label{eq:original-23}
\end{equation}
The final exponent \(\eta>0\) may be prescribed in advance and
made arbitrarily small.
\end{lemma}
\begin{proof}
\par\smallskip\noindent\textit{The surviving stub diagonal.}\quad
The single-intermediate exception of
Lemma~\ref{lem:early-index-words}, with both outers equal to \(P\),
is the copy increment
\[
 Z_x=h^2[p_x(u)^2e_v(v)+p_x(v)^2e_u(u)].
\]
There is a fixed \(C_4<q^2\) such that
\begin{equation}
 \sum_{x\in S,\ x\ne u}|p_x(u)|^4\le C_4.
\label{eq:original-24}
\end{equation}
Indeed, let \(s_u\) be the number of unresolved stubs at the vertex
of \(u\). Short paths of length less than \(g_0/3\) are unique
in total at each endpoint. The degree-zero contribution, excluding
\(u\), is at most \(s_u-1\). The positive-length contribution is
bounded by
\[
 \sum_{j\ge1}d(d-s_u)q^{j-1}q^{-2j}
 =\frac{d(d-s_u)}{q(q-1)}.
\]
For \(q=2\), the sum is at most 3; for \(q\ge3\), its maximum
over \(1\le s_u\le d\) is strictly below \(q^2\).
The long-path and uniform terms have vanishing fourth norm by
\eqref{eq:original-4}--\eqref{eq:original-7}. This proves
\eqref{eq:original-24} with a strict constant margin.

Summing over surviving \(x\) before averaging gives
\begin{equation}
 \mathbb E\sum_{x\ne u,v}Z_x^2
 \le(1+o(1))\frac{2h^4C_4}{l}T_S+Cd_m^2/l.
\label{eq:original-25}
\end{equation}
For its two squared summands use \eqref{eq:original-24} and then
sum the opposite diagonal square. For the cross product use
\[
 \sum_x\left(\sum_u p_x(u)^2\right)^2\le Cl,
 \qquad |e_u(u)e_v(v)|\le d_m^2.
\]
The conditional restriction has density \(1+o(1)\), so it preserves
the strict margin.

After subtracting \(Z_x\), the remainder obeys
\begin{equation}
 \mathbb E\sum_{x\ne u,v}|\Delta e_x(x)-Z_x|^2
 \le C(n^\eta/l+d_m^4).
\label{eq:original-26}
\end{equation}
Here is the complete reduction to the index-word cases.
A distinct-end word has second moment \(C/l^2\) for each \(x\),
by the first row of the averaging table. A non-pure same-end word
has second moment \(C/l^3\). A pure same-end word with topology
and at most one intermediate \(E\) has an opposite local factor
or a cross-topology factor; its second moment is
\(Cn^\eta/l^2\) by the third or fourth row of the table.
A word with at least two intermediate \(E\)'s has second moment
\(Cd_m^4/l\). If exactly one intermediate is \(E\), no topology
occurs, and an outer is also \(E\), use
\[
 \sum_s |e_x(s)|^2r_x(s)^2\le Cd_m^2
\]
in addition to the intermediate bound \(d_m\); this again gives
\(Cd_m^4/l\). The only same-end word left after these tests is
\(Z_x\), already removed. The formal tail follows the same tests,
with no \(E\)'s; its zero-intermediate pure-swap word is
necessarily distinct-end. The recentering terms contribute at most
\(Cd_m^2/l\) after summing over \(x\). Summing the per-output
bounds over at most \(l\) outputs and using the geometric
Minkowski bound in Lemma~\ref{lem:early-index-words} proves
\eqref{eq:original-26}.

Deletion removes \((1-o(1))2T_S/l\) in expectation. The existing
first-order square term is at most
\(2\mu_l\sqrt{lT_S}\). Since \(h^4C_4<1\), choose a fixed
\(\varepsilon>0\) such that
\((1+\varepsilon)h^4C_4<1\). Apply
\((x+y)^2\le(1+\varepsilon)x^2+(1+\varepsilon^{-1})y^2\)
to \(Z_x\) and its remainder and then Young's inequality to the
first-order term. Some fixed part of the deletion remains available
to absorb it. The resulting upper drift is
\[
 \mathbb E T_S'\le T_S+
 C(l^2\mu_l^2+n^\eta/l+d_m^4).
\]
The non-logarithmic quantities controlling the sum are
\[
 \frac{n^{2v_*}}{\alpha_*^2l_0}=n^{-.900a},\quad
 \frac{q^L}{l_0}=n^{-.580a+o(1)},\quad
 \frac{D_g^2}{l_0}=o(1),\quad
 nd_m^4=n^{-1/3+1.020a+o(1)}=o(1).
\]
The other term in \(\mu_l\) costs only \(n^{2\eta}\log n\).
A slightly larger small-power normalization makes the complete
additive budget summably small.

\par\smallskip\noindent\textit{Global diagonals.}\quad
No deletion is used. The copy squares have expected sum at most
\(CT_S/l+Cnd_m^2/l^2\), by the bounded fourth columns in
Lemma~\ref{lemma:pencil-coefficients}. The classification just proved gives
\(Cn(n^\eta/l^2+d_m^4/l)\) for the other squares.
The first-order term is at most \(2\mu_l\sqrt{nT_V}\).
At scale \(S_V=n^{\eta'}n/l_0\), the corresponding multiplier is
bounded by
\(C\mu_l\sqrt{l_0}/n^{\eta'/2}\), using
\(2\sqrt{T_V/S_V}\le1+T_V/S_V\).
The sum \(\sqrt{l_0}\sum\mu_l\) is at most an arbitrarily
small power, and the summed variance divided by \(n/l_0\)
has the same property. Choosing \(\eta'\) larger than the input
losses proves the global diagonal bound with an arbitrarily small
final loss.

\par\smallskip\noindent\textit{Existing near pairs.}\quad
For a fixed \(k\), there are \(O_k(l)\) near outputs. All
non-copy variance terms cost \(C_k(n^\eta/l+d_m^4)\).
For a copy term, at a fixed endpoint \(u\),
\[
 \sum_{i,x:\,\operatorname{dist}(i,x)\le k}
           A_i(u)^2A_x(u)^2\le C_k.
\]
This follows by summing first over the bounded ball around \(x\),
using the bounded sup of \(A_i(u)\), and then its squared column
sum over stub \(x\). Consequently the copy variance costs
\(C_kT_S/l\), rather than a self term involving \(T_{XS,k}\).
The existing first-order term is
\(C_k\mu_l\sqrt{lT_{XS,k}}\), and
\(\sum\sqrt l\mu_l\) is bounded by an arbitrarily small power.

\par\smallskip\noindent\textit{New near pairs.}\quad
A newly near pair has a short simple path using the new edge once.
For one of its two orientations, \(i\) is at old distance
\(k_1\) from \(u\), \(x\) at old distance \(k_2\) from \(v\),
and \(k_1+1+k_2\le k\). Each choice has bounded multiplicity.
The expected sum of the squared new entries is at most
\begin{equation}
 C_k\left(n^\eta+T_S+
       \sum_{X'\in\{V,S\},\ j<k}T_{X'S,j}\right)/l.
\label{eq:original-27}
\end{equation}
To verify this without discarding an aligned large reference entry,
expand the old entry plus every index word. An old entry, or an
outer \(E\) whose endpoint is opposite the associated short
neighborhood, is summed by its squared row bound over the free
opposite stub; bounded reverse neighborhood multiplicity gives
\(C_k/l\). An aligned outer \(e_i(u)\) or \(e_x(v)\) charges
\(T_{XS,k_1}/l\) or \(T_{SS,k_2}/l\), respectively. An
intermediate \(E\) charges \(T_S/l\) when diagonal and \(C/l\)
when off-diagonal, again by its squared rows. Other factors have
bounded sup and can be summed geometrically.

It remains to check a pure reference-tail word. Any intermediate
topology factor has second moment at most \(n^\eta/l\), even
with the bounded short-neighborhood multiplicities; a misaligned
outer has the same bound by its squared rows. For aligned outers,
split off their unique short monomials. Lemma~\ref{lem:early-near-reference}
charges the residual squares by
\(n^\eta+C(1+L)^4q^L/l\), and the free opposite endpoint supplies
\(1/l\). If both outers are their short monomials and there are
no intermediate factors, a pure swap has total degree
\(k_1+1+k_2<L\), so contributes nothing to the tail. A non-pure
swap has coefficient \(O(1/l)\), which is already within the
claimed bound. This exhausts the reference-tail cases. Recentring
costs a smaller term, since its sup is \(Cd_m/l\).
Minkowski's inequality and the geometric word summation prove
\eqref{eq:original-27}.

Finally choose the losses in a fixed finite order. Start with a
structural loss \(\eta_0>0\) so small that
\((20+4R_c)\eta_0<\eta\). Choose the normalization and extraction
exponents for \(T_S\), \(T_V\), and then successively for
\(T_{XS,0},\ldots,T_{XS,R_c}\), increasing the available exponent
by at least \(2\eta_0\) at each use of a preceding estimate.
Enlarge the first normalization enough to exceed twice the exponent
in its cumulative first-order budget. Every cumulative additive
cost is then \(o(1)\) after normalization, and every multiplier
has bounded sum. There are only \(2R_c+O(1)\) accounts per tracked
point. All extraction losses fit below the prescribed \(\eta\)
by taking \(\eta_0\) smaller if required. This proves
\eqref{eq:original-23}.
\end{proof}

\subsection{From averaged errors to uniform entry bounds}

The squared-sum estimates now supply a small factor in every mixed
moment. Pure powers require more care: their drift must have a
coefficient independent of the fixed moment order. We keep pair moments
and diagonal moments separate, then choose a fixed weight for the latter.
This absorbs the large transfer from pair moments to diagonal moments
without enlarging the common drift coefficient.

Set \(z_{ij}=E_{ij}/B_e\). For \(X,Y\in\{V,S\}\), let
\[
 \Phi_{XY}=\frac1{|X||Y|}\sum_{i\in X,j\in Y}|z_{ij}|^p,
 \qquad
 \Psi_X=\frac1{|X|}\sum_{i\in X}|z_{ii}|^p,
 \qquad F_p=\sum_{X,Y}\Phi_{XY},\quad D_p^{\rm mom}=\sum_X\Psi_X.
\]
Here \(p\ge4\) is an even integer, fixed independently of \(n\),
and \(M_p=1+F_p+D_p^{\rm mom}\). The notation
\(D_p^{\rm mom}\) is used only in this subsection and is distinct
from the two-by-two precision deviation \(D_p\).
All output averages are deterministic uniform sums in the current state.

By \eqref{eq:original-4}, \eqref{eq:original-7},
\eqref{eq:original-10}, and \eqref{eq:original-23}, some fixed
\(\beta>0\) has the following properties. Every row average of
\(|z|^2\), over either output type, is at most \(n^{-\beta}\).
Every diagonal second average has the same bound. For any fixed radius
\(R_c\), the measure
\[
 \int f\,d\mu_X=\frac1{|X|}
    \sum_{\substack{i\in X,u\in S\\
                    \operatorname{dist}(i,u)\le R_c}}f(i,u)
\]
has bounded mass and satisfies
\(\int|z_{iu}|^2d\mu_X\le n^{-\beta}\) and
\(\int|z_{iu}|^pd\mu_X\le l\Phi_{XS}\).
For example, before tiny losses the stub-row and near-pair second
averages have margin \(.62a\), whereas a global row has margin at
least \(.97a\). Thus \(\beta\) can be fixed, say below \(.30a\),
independently of the later choice of \(p\).

\begin{lemma}[Interpolation and a kernel inequality]\label{lem:early-interpolation}
Let a measure of bounded mass satisfy \(\int|z|^2\le n^{-\beta}\)
and \(\int|z|^p\le M\). For every fixed \(0<t<p\),
\begin{equation}
 \int|z|^t\le n^{-\beta_t}(1+M)^{t/p},
 \qquad \beta_t>0,
\label{eq:original-28}
\end{equation}
for sufficiently large \(n\). If a nonnegative kernel \(K\) on two
probability spaces has both marginal densities at most \(b\), then
for nonnegative \(f,g\) and \(r,s\ge0\), \(r+s\le1\),
\[
 \iint K(x,y)f(x)^rg(y)^s\,dx\,dy
 \le b\left(\int f\right)^r\left(\int g\right)^s.
\]
For \(K(i,u)=|p_i(u)|^k\), \(2\le k\le p\), with uniform
\(i\in X,u\in S\), one may use \(b=n^{\eta_0}/l\) for any
fixed \(\eta_0>0\), once \(n\) is sufficiently large. The row
marginal alone is at most \(C_p/l\). For \(k=p\) or \(2p\),
the constants in the row and column sums can be bounded by a constant
depending only on \(d\), for each fixed \(p\) and large \(n\).
\end{lemma}
\begin{proof}
For \(t\le2\), the second moment and H\"older's inequality give
\(C n^{-\beta t/2}\). For \(2<t<p\), interpolation gives
\[
 (n^{-\beta})^{(p-t)/(p-2)}M^{(t-2)/(p-2)}.
\]
The exponent of \(M\) is at most \(t/p\). Absorb fixed constants
by halving the positive exponent of \(n\), proving
\eqref{eq:original-28}.
For the kernel inequality, apply H\"older to the measure with density
\(K\): its mass, \(f\) integral, and \(g\) integral are at most
\(b\), \(b\int f\), and \(b\int g\).
For the stated kernel, its row sum is bounded by the squared row
bound and the bounded entry sup. The column sum for \(k=2\) is
polylogarithmic by \eqref{eq:original-7}; for \(k\ge3\) it is
bounded. Since \(|X|\ge c_dl\), this proves the marginal claim.
Finally Lemma~\ref{lemma:pencil-coefficients} bounds the relevant large-power column
sums uniformly in the fixed order, and the squared row bound gives
\(C_d(1+o(1))^{k-2}\) for row sums. With \(k\) fixed, this is
at most \(2C_d\) for large \(n\).
\end{proof}

\begin{lemma}[Linearization and its residual]\label{lem:early-linearization}
For every surviving output pair, decompose its actual increment as
\(\Delta E_{ij}=B_e\mathcal L_{ij}+\mathcal N_{ij}\), where
\(\mathcal L\) is linear in \(z\) and evaluated on the old output
vectors. The residual contains the nonlinear terms, the formal tail,
and the recentering of surviving stub outputs. It has sup
\[
 |\mathcal N_{ij}|\le C(d_m^2+B_{\rm tail}+d_m/l)
\]
and conditional second moment
\begin{equation}
 \mathbb E|\mathcal N_{ij}|^2
 \le C(n^\eta/l^2+d_m^4/l).
\label{eq:original-29}
\end{equation}
The leading terms of \(\mathcal L_{ij}\), up to their signs, are
\begin{equation}
 h p_i(u)z_{vj},\quad h p_i(v)z_{uj},\quad
 h p_j(u)z_{iv},\quad h p_j(v)z_{iu},
\label{eq:original-30}
\end{equation}
and
\begin{equation}
 h^2p_i(s)p_j(t)z_{\bar s\bar t},
 \qquad s,t\in\{u,v\}.
\label{eq:original-31}
\end{equation}
All other coefficients have sup bounded by a fixed negative power
of \(n\). A diagonal coefficient of \(H_P\) satisfies, with its
first endpoint \(s\) fixed,
\begin{equation}
 \mathbb E_{\bar s}\bigl[|(H_P)_{ss}|^k
                  \mathbf 1_{\rm valid}\bigr]
 \le C_p n^\eta n^{-c'_0(k-2)}/l,
 \qquad 2\le k\le p,
\label{eq:original-32}
\end{equation}
where the expectation may also omit boundedly many output labels.
\end{lemma}
\begin{proof}
Put \(T_P=I-PUH_PU^T\). The exact inverse-update difference identity
is
\[
 (W-\mathcal Q(W))-(P-\mathcal Q(P))
 =T_P(E-EUH_WU^TE)T_P^T.
\]
It follows either by multiplying the two inverse-update formulas or
by using
\(H_W-H_P=-H_P(U^TEU)H_W\).
Thus the linear old-axis increment is \(T_PET_P^T-E\).
Replacing \(H_P\) by \(-\sigma hF\) gives
\eqref{eq:original-30} and \eqref{eq:original-31}.
The difference \(H_P+\sigma hF\) is power-small by
\eqref{eq:original-5} and its geometric series. The leftover
one-axis forms are \(p_i(s)z_{tj}\) and their transposes; the
two-axis forms are \(p_i(s)p_j(t)z_{cd}\), with coefficient
\((H_P)_{sc}(H_P)_{dt}\) minus its pure-swap value.

The nonlinear term is \(-T_PEUH_WU^TET_P^T\), whose entries
have sup \(Cd_m^2\). Its complete second-moment classification is
as follows. Distinct outer endpoints use two squared rows and cost
\(C/l^2\). Same-end words with two intermediate \(E\) factors
cost \(Cd_m^4/l\). If one intermediate and one outer are \(E\),
use \(\sum_s|e_i(s)|^2r_j(s)^2\le Cd_m^2\) and the second
\(E\) bound to get the same estimate. If both \(E\)'s are
outer factors and no intermediate \(E\) occurs, a same-end word
has either a non-pure swap or an opposite local/cross-topology
factor; its second moment is at most \(Cn^\eta/l^2\) by the
averaging table. These are all possible placements of at least two
\(E\)'s in Lemma~\ref{lem:early-index-words}.
For a tail word there are no \(E\)'s: distinct-end words cost
\(C/l^2\), and same-end words have the non-pure or opposite-topology
factor and cost \(Cn^\eta/l^2\). Recentring has squared sup
\(Cd_m^2/l^2\). The geometric Minkowski bound now proves
\eqref{eq:original-29}; the sup claim follows also from
\eqref{eq:original-6}.

For \eqref{eq:original-32}, follow the first pure swap in the
\((s,s)\) entry of \(H_P\). By
Lemma~\ref{lem:early-index-words}, either a non-pure swap occurs,
or a local factor at \(\bar s\), or a cross factor occurs. Hence
\[
 |(H_P)_{ss}|\le C(1/l+j_{\bar s}+b_{s\bar s}).
\]
Its conditional second moment is \(Cn^\eta/l\); its sup is
\(n^{-c'_0}\). Multiply the second-moment bound by the
\((k-2)\)-th power of this sup. Omitting a bounded number of
partners only changes the bounded density constant.
\end{proof}

\begin{lemma}[High moments]\label{lem:early-high-moments}
There is a constant \(C_*\), independent of every sufficiently
large even order \(p\), with the following property. After \(p\)
is fixed, one may choose a fixed \(\kappa_p>0\), a fixed near
radius \(R_c\), and sufficiently small structural and low-square
losses such that
\begin{equation}
 Z_p=1+F_p+\kappa_pD_p^{\rm mom},\qquad
 \mathbb E Z_p'\le(1+C_*/l+u_l)Z_p,
 \qquad \sum_lu_l=o(1),
\label{eq:original-33}
\end{equation}
with nonnegative deterministic upper budgets \(u_l\) under the
bootstrap assumptions.
\end{lemma}
\begin{proof}
We separate mixed powers from pure powers. The low-square estimates and
interpolation make every mixed contribution a power-small multiple of
\(M_p/l\), whose sum is small. Pure powers can transfer pair moments
to diagonal moments at rate \(C_pF_p/l\); their reverse transfer has
the smaller rate \(C_pD_p^{\rm mom}/l^2\). Weighting the diagonal
family by \(\kappa_p\) makes the first coefficient bounded when
\(\kappa_p C_p\le1\), while the reverse coefficient
\(C_p/(\kappa_p l^2)\) remains summable for fixed \(p\).
Thus we must bound each family's own \(1/l\) drift independently of
\(p\), and verify these two transfer rates. A fixed threshold on the
outer path entries will distinguish the pure powers that need this
uniform bound from terms that gain a small factor.
All expectations of nonnegative quantities below may be taken over
independent uniform endpoints, losing a bounded density factor.
For signed first-order terms we retain the survival-conditional
mean \eqref{eq:original-22}. Throughout this proof constants marked
\(C_p\) may depend on \(p\), whereas \(C_0,C_1\) do not.

\par\smallskip\noindent\textit{First-order and residual terms.}\quad
In the even-power Taylor formula, the first-order term has absolute
mean at most \(p(\mu_l/B_e)|z_{ij}|^{p-1}\).
The sum \(\sum_l\mu_l/B_e\) tends to zero by a power.
Write \(r_{ij}=\mathcal N_{ij}/B_e\). Its sup is a fixed
negative power of \(n\). For \(2\le k\le p\),
\[
 |z_{ij}|^{p-k}\mathbb E|r_{ij}|^k
 \le C(1+|z_{ij}|^p)
       \frac{n^\eta/l^2+d_m^4/l}{B_e^2}.
\]
The sum of the displayed coefficient tends to zero.
For \(k<p\), finite-term splitting costs only \(C_p\).
For \(k=p\), use
\(|a+b|^p\le2|a|^p+C_p|b|^p\) when separating the retained
linear part from the residual. This incurs a factor 2, independent
of \(p\), on the retained part.

\par\smallskip\noindent\textit{Intermediate powers of linear terms.}\quad
Fix \(2\le k<p\). The minimum of the positive interpolation
exponents in \eqref{eq:original-28}, for the finitely many powers
used below, is positive. Choose the kernel loss \(\eta_0\) less
than half this minimum.
For the first term in \eqref{eq:original-30}, integrating \(u\)
uses \(\sum_u|p_i(u)|^k\le C_p\), producing \(C_p/l\).
For fixed \(j\), interpolation in the free \(v\) row gives
\[
 \mathbb E_v|z_{vj}|^k
 \le n^{-\delta}(1+R_j)^{k/p},\qquad
 R_j=\mathbb E_v|z_{vj}|^p.
\]
H\"older in the remaining output indices bounds the result by
\(C_p n^{-\delta}M_p/l\). This derivation covers each of the four
orientations and also diagonal outputs \(j=i\).

For \eqref{eq:original-31} with \(s=t\), discard one bounded
outer factor. Integration of the other outer over \(s\) gives
\(C_p/l\), and interpolation in the independent diagonal input
\(z_{\bar s\bar s}\) gives the power saving. If \(s\ne t\),
retain \(|p_i(s)|^k\) and discard the other bounded outer.
Interpolation over \(t\), conditional on \(s\), yields
\(n^{-\delta}(1+R_s)^{k/p}\). Average the output factor first
over \(j\), and apply the kernel inequality in
Lemma~\ref{lem:early-interpolation} to \((i,s)\), with exponents
\((p-k)/p\) and \(k/p\). The cost is
\(C_p n^{\eta_0-\delta}M_p/l\). For a diagonal output, use
\(|z_{ii}|^p\) instead of its output-row average in that kernel
inequality. This proves the same bound.

A leftover one-axis term whose input endpoint differs from its
outer endpoint uses the preceding one-axis calculation. If these
endpoints coincide, its coefficient is \((H_P)_{ss}\); integrate
the other selected endpoint first by \eqref{eq:original-32}, then
use the row interpolation in \(z_{sj}\). Here the structural loss
\(\eta\) is chosen below the interpolation saving, including when
\(k=2\) in \eqref{eq:original-32}. For a leftover two-axis
term retain the outer at \(s\). If its input is \(z_{ss}\),
the coefficient includes \((H_P)_{ss}\), so
\eqref{eq:original-32} and diagonal interpolation apply after
discarding the bounded outers. An input \(z_{\bar s\bar s}\)
uses independent diagonal interpolation; an off-diagonal input uses
the kernel argument just given. These three input types exhaust
\(\{u,v\}^2\). Thus every intermediate Taylor power costs
\(C_p n^{-\delta_p}M_p/l\) for some \(\delta_p>0\).

\par\smallskip\noindent\textit{Pure powers for independent output pairs.}\quad
The two-axis pure powers can be charged at once:
\[
 \mathbb E_{i,j,u,v}
   |p_i(s)p_j(t)z_{cd}|^p
 \le \frac{C_p}{|X||Y|}\mathbb E_{u,v}|z_{cd}|^p
 \le C_pM_p/l^2.
\]
This holds also for diagonal inputs \(c=d\), because the outputs
\(i,j\) remain independent. The leftover one-axis pure powers
have either the pointwise power-small coefficient and a row
integration, or \eqref{eq:original-32} when input and outer
endpoints coincide. Hence they cost a power-small multiple of
\(M_p/l\).

It remains to sum the leading one-axis terms without multiplying
all their pure powers by \(C_p\). Choose a fixed threshold
\(\varepsilon_p>0\). The sum of terms with
\(|p_i(u)|\le\varepsilon_p\) has pure moment at most
\[
 C_p\varepsilon_p^{p-2}F_p/l.
\]
This follows from
\(|p_i(u)|^p\mathbf1_{|p_i(u)|\le\varepsilon_p}
 \le\varepsilon_p^{p-2}|p_i(u)|^2\)
and the row-sum bound. The threshold will be chosen last to make
this coefficient as small as required.
Every outer above the threshold is at actual distance at most a
fixed \(R_c=R_c(p,d)\) from its selected endpoint, by the
short-path and long-entry estimates in
Lemma~\ref{lemma:pencil-coefficients}.
For sufficiently large \(n\), \(g_0>2R_c\), so an output axis
has at most one such endpoint.

If both output axes have large outers at the same endpoint, say
\(u\), the corresponding pure powers have an extra \(1/l\)
gain even after a \(C_p\) splitting. To see it for
\(p_i(u)z_{vj}\), sum \(|p_i(u)|^p\) over \(i\), using its
bounded column sum. Each \(j\) has only \(C_{R_c,d}\) stubs
\(u\) in its near ball. Therefore
\[
 \mathbb E_{i,j,u,v}
 \bigl[|p_i(u)z_{vj}|^p\mathbf1_{j\text{ near }u}\bigr]
 \le C_p\Phi_{SY}/(|X|l)\le C_pF_p/l^2.
\]
This calculation remains valid when several unresolved stubs have
the same vertex; there are at most \(d\) such labels.

If the two large outers occur at opposite endpoints, there are at
most two retained terms. Expand their \(p\)-th power, keeping
both pure summands with coefficient 1. Each mixed term has positive
powers \(t,p-t\), with \(0<t<p\), and is at most
\[
 \frac{C_p}{l^2}
  \left(\int|z_{iu}|^{p-t}\,d\mu_X\right)
  \left(\int|z_{jv}|^t\,d\mu_Y\right).
\]
By \eqref{eq:original-28} and the near-pair bounds, this is at most
\(C_p n^{-\delta_p}(1+F_p)/l\). A single large term has no
mixed powers. Finally the pure powers of all retained one-axis
terms satisfy
\[
 \mathbb E\sum_{\rm orientations}|h p_i(u)z_{vj}|^p
 \le C_0h^pF_p/l.
\]
Here \(C_0\) is independent of \(p\): there are only four
orientations and four output types, their squared row constants
depend only on \(d\), and their row \(p\)-sums are bounded by
\(C_d(1+o(1))^{p-2}\le2C_d\) for each fixed \(p\) and large
\(n\). The Young splittings separate a number of groups independent
of \(p\) and retain the principal part with only fixed relative
factors (at most 2 per splitting),
and conditioning costs a bounded density factor, not a \(p\)-th
power. Thus these operations only enlarge \(C_0\) by a constant
independent of \(p\).

\par\smallskip\noindent\textit{Pure powers for diagonal outputs.}\quad
For a diagonal \(i=j\), an above-threshold orientation, say at
\(u\), has retained terms
\[
 A=2\sigma h p_i(u)z_{vi},\qquad
 B=h^2p_i(u)^2z_{vv}.
\]
Their pure powers obey
\[
 \mathbb E|A|^p\le C_pF_p/l,
 \qquad \mathbb E|B|^p\le C_0h^{2p}D_p^{\rm mom}/l.
\]
The first is explicitly a pair-to-diagonal transfer. The second
uses a row \(2p\)-sum with a constant independent of the fixed
order, as above. For a mixed power \(|A|^t|B|^{p-t}\), first
sum over \(u\); the outer power is \(2p-t\ge p+1\), so that
sum contributes \(C_p/l\). Interpolate over \(i\) at fixed
\(v\) in the positive power \(|z_{vi}|^t\), and then apply
H\"older in \(v\) against \(|z_{vv}|^{p-t}\). This gives
\(C_p n^{-\delta_p}M_p/l\).

The other diagonal two-axis term has off-diagonal input \(z_{uv}\).
At least one of its two outer factors is below the threshold.
Keep the other column \(p\)-power and sum it over \(i\); this
costs \(C_p\varepsilon_p^p\Phi_{SS}/l\).
A below-threshold same-end diagonal input costs
\(C_p\varepsilon_p^{2p-2}D_p^{\rm mom}/l\).
Below-threshold one-axis terms cost
\(C_p\varepsilon_p^{p-2}F_p/l\).
For a leftover diagonal one-axis term
\(a_{st}p_i(s)z_{ti}\), distinguish its input endpoint. If
\(t=s\), its coefficient is \(a_{ss}=(H_P)_{ss}\).
Average the opposite selected endpoint first using
\eqref{eq:original-32}, discard the bounded outer factor, and
then average \((i,s)\). This gives
\[
 \mathbb E|a_{ss}p_i(s)z_{si}|^p
 \le\frac{C_p n^{\eta-c'_0(p-2)}}l\Phi_{SX}.
\]
If \(t=\bar s\), the leftover coefficient has sup
\(n^{-c}\) for some fixed \(c>0\). At fixed \(i\), sum the
row \(p\)-power \(|p_i(s)|^p\) over \(s\), obtaining
\[
 \mathbb E|a_{s\bar s}p_i(s)z_{\bar s i}|^p
 \le\frac{C_p n^{-cp}}l\Phi_{SX}.
\]
For a leftover two-axis term, the input \(z_{cd}\) is independent
of the output \(i\). Discard one bounded outer factor and sum
the other column \(p\)-power over \(i\), which gives
\(C_p n^{-cp}(F_p+D_p^{\rm mom})/l\).
For sufficiently large fixed \(p\) and sufficiently small losses,
all these bounds are \(C_p n^{-\delta_p}M_p/l\).
Thus every diagonal pure-power term has now been accounted for.

\par\smallskip\noindent\textit{Combining pair and diagonal moments.}\quad
Collecting the self-drift, the two transfers between moment families, and
the summable errors gives, for a fixed \(C_1\),
\[
\begin{split}
 \mathbb E F_p'&\le
   \left(1+\frac{C_1+C_0h^p+\epsilon}{l}\right)F_p
   +\frac{C_p}{l^2}D_p^{\rm mom}+v_lM_p,\\
 \mathbb E(D_p^{\rm mom})'&\le
   \left(1+\frac{C_1+C_0h^{2p}+\epsilon}{l}\right)D_p^{\rm mom}
   +\frac{C_p}{l}F_p+v_lM_p,
\end{split}
\]
where \(\epsilon>0\) is any prescribed constant after the threshold
is chosen, and \(\sum_l v_l=o(1)\). To obtain \(v_l\), sum the
first-order and residual budgets and the terms
\(C_pn^{-\delta_p}/l\); all have summably small totals.
The coefficient of the nonsummable transfer
\(C_pF_p/l\) depends only on \(p,d\) and the uniform row bounds.
Constants depending on the resulting near radius \(R_c\) occur
only with \(l^{-2}\) or with the power saving
\(n^{-\delta_p}/l\), and are included in \(v_l\).
The change of a denominator \(|S|\) to \(|S|-2\) contributes
at most \(1+C_1/l\), with \(C_1\) independent of \(p\).
Deletion of labels only lowers their unnormalized moment sums.

Choose \(p_0\) so large that \(C_0q^{-p_0/2}\le1\).
For each even \(p\ge p_0\), choose
\(0<\kappa_p\le\min(1,(1+C_p)^{-1})\), then choose the
threshold so that \(\epsilon\le1\).
Since \(h\le q^{-1/2}\), multiplication of the second inequality
by \(\kappa_p\) and addition give
\eqref{eq:original-33} with, for example,
\(C_*=C_1+4\). The reverse leakage
\(C_pD_p^{\rm mom}/l^2\) costs
\(C_p/(\kappa_pl^2)\) in the multiplier and has sum \(o(1)\).
The multiplier from \(v_lM_p\) is at most a fixed
\(C_p\kappa_p^{-1}v_l\), whose sum also tends to zero.
This proves the required independence of \(C_*\) from \(p\).
\end{proof}

\begin{proof}[Proof of Proposition~\ref{prop:early-phase}]
At the empty state the accounts have the asserted initial values.
Assume their conclusions at the current state. Each stub has at
most \(D_g=O_d(q^{g_0})=o(l_0)\) forbidden partners, so
\(l-1-D_g>0\) throughout this phase and the candidate set is
nonempty. The deterministic
implications \eqref{eq:original-8}--\eqref{eq:original-12} give the
row, trace, and operator estimates used in the transition lemmas.
Every available pair is small in its two-by-two deviation, so its
next precision is positive by the inverse-update inertia argument.
The structural accounts, Lemma~\ref{lem:early-trace-accounts},
Lemma~\ref{lem:early-low-squares}, and
Lemma~\ref{lem:early-high-moments} can therefore be included in the
same discounted total.

The high-moment discount loss is at most \(n^{C_*+1}\) for large
\(n\) since \(\sum_l l^{-1}=O(\log n)\). Extraction of one entry
from a uniform pair average costs at most \(d^2n^2\), and from a
diagonal average costs less. Fixed \(\kappa_p^{-1}\) is harmless.
Choose the even order \(p\) large enough that
\((C_*+4)/p<.02a\). Then the discounted total and
\eqref{eq:original-33} imply
\(\max|E_{ij}|<B_en^{.02a}=d_m\), with a positive power of slack.
The trace and low-square accounts have their own exponent cushions.
Thus account extraction implies the required bounds at the chosen
next state, closing the induction.

There is no reverse dependence in these choices. First fix a
sufficiently large even \(p\) using the order-independent \(C_*\);
then choose \(\kappa_p\), the threshold and its finite \(R_c\);
then decrease the structural and low-square losses for the finitely
many interpolations and near-radius inductions. Their constants can
grow but remain independent of \(n\). This proves the proposition.
\end{proof}
